\providecommand{\kBKM}{\kappa_{\mathrm{BKM}}}
\providecommand{\kch}{\kappa_{\mathrm{ch}}}
\providecommand{\kcat}{\kappa_{\mathrm{cat}}}
\providecommand{\kHeis}{\kappa_{\mathrm{ch}}^{\mathrm{Heis}}}
\providecommand{\kfib}{\kappa_{\mathrm{fiber}}}
\providecommand{\Z}{\mathbb{Z}}
\providecommand{\R}{\mathbb{R}}
\providecommand{\cO}{\mathcal{O}}
\providecommand{\Hup}{\mathbb{H}_{2}}
\providecommand{\wt}{\operatorname{wt}}
\providecommand{\DeltaCHLThree}{\Delta_{3}^{(3)}}
\providecommand{\DeltaDDThree}{\Delta_{1,3}}
\providecommand{\NablaThree}{\nabla_{2}}
\providecommand{\GammaDDThree}{\Gamma_{3}}
\providecommand{\GammaCHLThree}{\Gamma_{0}^{+}(3)}
\providecommand{\gDeltaCHLThree}{\mathfrak{g}_{\Delta_{3}^{(3)}}}
\providecommand{\fgGNThree}{\mathfrak{g}_{3}^{\mathrm{GN}}}
\providecommand{\varphiCHLThree}{\varphi_{3}^{\mathrm{CHL}}}
\providecommand{\phiDDThree}{\phi_{0,3}}
\providecommand{\GammaDDProd}{\Gamma_{3}^{\mathrm{prod}}}
\providecommand{\Hdiag}{\mathcal{H}_{\mathrm{diag}}}

\section*{Order-three automorphic products}
\label{sec:N3-automorphic-products}

\paragraph{Two products.}
Two genus-two automorphic products occupy the order-three boundary.
The primitive CHL denominator is
\[
  \DeltaCHLThree\in M_3(\GammaCHLThree,\chi_3),
\]
attached to a symplectic K3 automorphism \(g_3\) of Mathieu class
\(3A\), frame shape \(1^6 3^6\), and primitive CHL Jacobi input
\(\varphiCHLThree\).  The diagonal-divisor product in the
Gritsenko-Cl\'ery catalogue is
\[
  \DeltaDDThree\in M_1(\GammaDDThree,\chi_6),
  \qquad (t,N_{\mathrm{lvl}};k)=(3,1;1).
\]
This is the product denoted \(\Delta_1\) in the
Gritsenko-Cl\'ery row order and \(\Delta_{1}^{(3)}\) in the
level-three denominator notation.
Its Jacobi input is the Gritsenko-Nikulin form
\[
  \phiDDThree(\tau,z)
  =
  \left(\frac{\vartheta(\tau,2z)}{\vartheta(\tau,z)}\right)^2,
\]
and its product lattice is
\[
  M_3=U(12)\oplus\langle2\rangle.
\]
The second catalogue row at level \(3\) is
\[
  \NablaThree\in M_2(\Gamma^{(2)}_0(3),\chi_2),
  \qquad (t,N_{\mathrm{lvl}};k)=(1,3;2).
\]
The CHL order \(3\) and the catalogue coordinate \(t=3\) label
distinct products.  A row-map theorem requires a common host, chamber,
multiplier, and coefficient convention.  The automorphic sentence
contains the product, group, character, divisor, lattice, and Weyl
chamber.  The positive-half Hall construction, the Hall-Drinfeld
recognition problem, and the scalar DT convention are separate data.

\paragraph{Primitive CHL construction.}
Let \(E\) be an elliptic curve and let \(t_3\in E[3]\) be a primitive
translation.  For \(a=1,2\), the translation \(t_3^a\) acts freely on
\(E\); hence \((g_3^a,t_3^a)\) acts freely on \(K3\times E\).  The
quotient
\[
  X_3=(K3\times E)/\langle(g_3,t_3)\rangle
\]
is a smooth Calabi--Yau threefold.  The cover \(K3\times E\to X_3\) is
finite etale of degree \(3\).  The action preserves the product
holomorphic volume form, and the quotient of the trivial canonical
bundle is trivial.

\paragraph{Primitive CHL weight.}
Let
\[
  \varphiCHLThree(\tau,z)
  =
  \sum_{n,\ell}c^{\mathrm{CHL}}_3(n,\ell)q^n\zeta^\ell
\]
be the primitive order-three CHL Borcherds input of weight \(0\) and
index \(1\), in the denominator normalisation.  Its zeroth coefficient is
\[
  c_3(0)=c^{\mathrm{CHL}}_3(0,0)=6.
\]
The primitive CHL table is
\[
  (c_N(0))_{N=(1,2,3,4,6)}=(10,8,6,4,2).
\]
Borcherds' singular theta correspondence gives
\[
  \wt(\DeltaCHLThree)
  =\frac{c_3(0)}{2}=3,
  \qquad
  \kBKM(\DeltaCHLThree)=3.
\]
Equivalently,
\[
  (\kBKM(\Phi_N))_{N=(1,2,3,4,6)}=(5,4,3,2,1).
\]
The equality \(c_3(0)=6\) also matches the Lefschetz fixed-point count
of \(g_3\) on the K3 fibre.  The definition used in the Borcherds
weight theorem is the Fourier coefficient above.
The scalar dyon convention uses \((\DeltaCHLThree)^2\) and has weight
\(6\) at \(N=3\).  The BKM denominator convention uses the primitive
form \(\DeltaCHLThree\) and has weight \(3\).

\paragraph{Primitive CHL product.}
For \(Z=(\tau,z,\sigma)\in\Hup\) and
\[
  (q,\zeta,p)
  =
  (e^{2\pi i\tau},e^{2\pi iz},e^{2\pi i\sigma}),
\]
the CHL product has the form
\[
  \DeltaCHLThree(Z)
  =
  q^A\zeta^B p^C
  \prod_{\substack{(n,\ell,m)\in\Z^3\\(n,\ell,m)>0}}
  \left(1-q^n\zeta^\ell p^m\right)^{
    c^{\mathrm{CHL}}_3(nm,\ell)} .
\]
The Weyl vector supplies the initial exponent \(q^A\zeta^B p^C\).
The coefficient \(c_3(0)=6\) fixes the weight through
Borcherds' formula; the remaining Fourier coefficients
\(c^{\mathrm{CHL}}_3(nm,\ell)\) are the product exponents.

\paragraph{Diagonal-divisor product.}
Gritsenko-Cl\'ery classify the genus-two Siegel modular forms with
diagonal divisor of multiplicity one for Hecke type congruence
subgroups \(\Gamma_t(N_{\mathrm{lvl}})\).  Their eight triples are
\[
\begin{gathered}
 (1,1;5),\ (2,1;2),\ (1,2;3),\ (3,1;1),\\
 (1,3;2),\ (4,1;\tfrac12),\ (1,4;\tfrac32),\ (2,2;1).
\end{gathered}
\]
The twice-weight tuple is
\[
  (10,4,6,2,4,1,3,2).
\]
Equivalently the weight tuple in triple order is
\[
  (5,2,3,1,2,\tfrac12,\tfrac32,1).
\]
The row \((3,1;1)\) has
\[
  \wt(\DeltaDDThree)=1,
  \qquad
  c^{\mathrm{dd}}_{3,1}(0,0)=2,
  \qquad
  \kBKM(\DeltaDDThree)=1.
\]
The divisor is
\[
  \operatorname{supp}\operatorname{div}_{\Hup}(\DeltaDDThree)
  =
  \bigcup_{\gamma\in\GammaDDThree}\gamma\langle\Hdiag\rangle,
  \qquad
  \operatorname{mult}_{\Hdiag}(\DeltaDDThree)=1.
\]
In the Gritsenko-Nikulin exponential normalisation set
\[
  M_3=\Z f_2\oplus\Z\widehat f_3\oplus\Z f_{-2},
  \qquad
  \left((f_i,f_j)\right)_{(f_2,\widehat f_3,f_{-2})}
  =
  \begin{pmatrix}
    0&0&-12\\
    0&2&0\\
    -12&0&0
  \end{pmatrix},
\]
and
\[
  \alpha_3(n,\ell,m)=nf_2+\ell\widehat f_3+mf_{-2},
  \qquad
  e_3^{(n,\ell,m)}=q^n\zeta^\ell p^{3m}.
\]
The product chamber \(\GammaDDProd\) is the Gritsenko-Nikulin positive
cone: \(m>0\), or \(m=0,n>0\), or \(m=n=0,\ell<0\).  With
\(\phiDDThree=\sum f_3(n,\ell)q^n\zeta^\ell\), the diagonal-divisor
product is
\[
  \DeltaDDThree(Z)
  =
  q^{1/6}\zeta^{1/2}p^{1/2}
  \prod_{(n,\ell,m)\in\GammaDDProd}
  \left(1-q^n\zeta^\ell p^{3m}\right)^{f_3(nm,\ell)}.
\]
The row \((1,3;2)\) has \(\wt(\NablaThree)=2\) and
\(c^{\mathrm{dd}}_{1,3}(0,0)=4\).  The constants
\(c^{\mathrm{dd}}_{3,1}(0,0)=2\) and \(c_3(0)=6\) belong to different
normalisations.

\paragraph{Denominator algebra.}
The row \(\DeltaDDThree\) carries the Gritsenko-Nikulin denominator
algebra \(\fgGNThree\) on the Lorentzian lattice \(M_3\).  Its Weyl
vector is
\[
  \rho_3=\left(\frac16,\frac12,\frac16\right),
\]
and its real simple roots are the \(P_{3,II,\overline0}\)-walls of the
Gritsenko-Nikulin chamber in \(M_3\otimes\R\).  The denominator identity
has signed root multiplicities \(f_3(nm,\ell)\).

A denominator theorem for the CHL product \(\DeltaCHLThree\) uses its
own Lorentzian root lattice, real simple roots, Weyl chamber, parity
split of the imaginary simple roots, and Weyl-Kac-Borcherds identity
with signed multiplicities \(c^{\mathrm{CHL}}_3(n,\ell)\).  A compact
Hall or CoHA realisation adds orientation, finite-stabiliser
linearisation, an integration map, and coefficient comparison with the
compact BPS moduli problem.

\paragraph{\(K3\times E\) invariants.}
The compact total-space invariants of \(K3\times E\) are
\[
  \kcat(K3\times E)
  =\chi(\cO_{K3})\chi(\cO_E)
  =2\cdot 0=0,
  \qquad
  \kch(K3\times E)
  =
  \sum_q(-1)^qh^{0,q}(K3\times E)
  =0.
\]
The Heisenberg specialisation and the untwisted BKM denominator weight
are
\[
  \kHeis(K3\times E)=3,
  \qquad
  \kBKM(\Delta_5)=\frac{c_1(0)}{2}=\frac{10}{2}=5.
\]
The fibre invariant is the Mukai lattice rank
\[
  \kfib(K3)=\operatorname{rk}\widetilde H(K3,\Z)=1+22+1=24.
\]
The labelled \(K3\times E\) constants are
\[
  \{\kcat,\kch,\kHeis,\kBKM,\kfib\}
  =
  \{0,0,3,5,24\}.
\]
The compact chiral supertrace is \(\kch(K3\times E)=0\).  The value
\(3\) belongs to the Heisenberg-Mukai specialisation.  The value \(5\)
belongs to the untwisted denominator \(\Delta_5\).  The value \(24\)
belongs to the K3 fibre.

\paragraph{\(N=1\) normalisation.}
The monic Borcherds product is
\[
  D_5=64^{-1}\Delta_5,
  \qquad
  \kBKM(\Delta_5)=5.
\]
The scalar normalisation changes neither the BKM weight nor the root
data.  The Oberdieck-Pandharipande reduced scalar uses
\[
  \Phi_{10}^{\mathrm{OP}}
  =
  D_5^2
  =
  4096^{-1}\Delta_5^2,
  \qquad
  Z_{\mathrm{OP}}
  =
  -(\Phi_{10}^{\mathrm{OP}})^{-1}
  =
  -4096\,\Delta_5^{-2}.
\]
The primitive denominator algebra is governed by the unsquared
Borcherds weight.

\paragraph{Hall comparison.}
The primitive CHL target is \(\DeltaCHLThree\).  The diagonal-divisor
targets \(\DeltaDDThree\) and \(\NablaThree\) live on their own
\((t,N_{\mathrm{lvl}})\) rows.  A BPS comparison on \(X_3\) consists of
equality of product exponents, compatibility with the
Kontsevich-Soibelman wall-crossing product, and compatibility with the
Hall bracket on the chosen compact or orbifold BPS category.  Since
\(X_3\) is a Calabi--Yau threefold, the native positive half is an
\(E_1\) Hall algebra.  Braided \(E_2\) structure belongs to the
Drinfeld centre, the Drinfeld double, or the representation/module
category built from that \(E_1\) algebra.

\paragraph{References.}
Borcherds' weight formula is Borcherds 1998, Theorem 13.3.
The primitive CHL constants are the Gritsenko-Nikulin and Gritsenko
denominator constants: Gritsenko-Nikulin 1995, Part II, Theorem 2.1,
and Gritsenko 1999, Corollary 3.3.  The diagonal-divisor row list and
its weights are Gritsenko-Cl\'ery, \emph{The Siegel modular forms of
genus two with the simplest divisor}, Theorems 1.4 and 3.1.  The
Gritsenko-Nikulin row \(\DeltaDDThree\) product and the input
\(\phiDDThree=(\vartheta(\tau,2z)/\vartheta(\tau,z))^2\) are
Gritsenko-Nikulin 1995, Part II, Theorem 2.1 and Section 5.  The Mukai
lattice rank is Mukai 1987, Section 1, with the standard K3 cohomology
presentation reviewed in Huybrechts 2016, Chapter 16.
Oberdieck-Pandharipande fix the reduced \(K3\times E\) scalar
normalisation.  Davison-Meinhardt and Kontsevich-Soibelman supply the
Hall-side comparison formalism after the orientation and purity data
are fixed.
